% Options for packages loaded elsewhere
\PassOptionsToPackage{unicode}{hyperref}
\PassOptionsToPackage{hyphens}{url}
\PassOptionsToPackage{dvipsnames,svgnames,x11names}{xcolor}
%
\documentclass[
  11pt,
]{article}
\usepackage{amsmath,amssymb}
\usepackage{iftex}
\ifPDFTeX
  \usepackage[T1]{fontenc}
  \usepackage[utf8]{inputenc}
  \usepackage{textcomp} % provide euro and other symbols
\else % if luatex or xetex
  \usepackage{unicode-math} % this also loads fontspec
  \defaultfontfeatures{Scale=MatchLowercase}
  \defaultfontfeatures[\rmfamily]{Ligatures=TeX,Scale=1}
\fi
\usepackage[]{lmodern}
\ifPDFTeX\else
  % xetex/luatex font selection
\fi
% Use upquote if available, for straight quotes in verbatim environments
\IfFileExists{upquote.sty}{\usepackage{upquote}}{}
\IfFileExists{microtype.sty}{% use microtype if available
  \usepackage[]{microtype}
  \UseMicrotypeSet[protrusion]{basicmath} % disable protrusion for tt fonts
}{}
\makeatletter
\@ifundefined{KOMAClassName}{% if non-KOMA class
  \IfFileExists{parskip.sty}{%
    \usepackage{parskip}
  }{% else
    \setlength{\parindent}{0pt}
    \setlength{\parskip}{6pt plus 2pt minus 1pt}}
}{% if KOMA class
  \KOMAoptions{parskip=half}}
\makeatother
\usepackage{xcolor}
\usepackage[margin=0.85in]{geometry}
\usepackage{longtable,booktabs,array}
\usepackage{calc} % for calculating minipage widths
% Correct order of tables after \paragraph or \subparagraph
\usepackage{etoolbox}
\makeatletter
\patchcmd\longtable{\par}{\if@noskipsec\mbox{}\fi\par}{}{}
\makeatother
% Allow footnotes in longtable head/foot
\IfFileExists{footnotehyper.sty}{\usepackage{footnotehyper}}{\usepackage{footnote}}
\makesavenoteenv{longtable}
\setlength{\emergencystretch}{3em} % prevent overfull lines
\providecommand{\tightlist}{%
  \setlength{\itemsep}{0pt}\setlength{\parskip}{0pt}}
\setcounter{secnumdepth}{-\maxdimen} % remove section numbering
\usepackage{amsmath,amssymb,mathtools}
\usepackage{microtype}
\usepackage{longtable,booktabs,array,etoolbox}
\usepackage{fancyhdr}
\usepackage{fvextra}
\usepackage{xurl}
\usepackage{titlesec}
\usepackage{needspace}
\definecolor{Ink}{HTML}{182D40}
\definecolor{Accent}{HTML}{176C78}
\definecolor{Muted}{HTML}{5D6871}
\titleformat{\section}{\Large\bfseries\color{Ink}}{}{0pt}{}
\titleformat{\subsection}{\large\bfseries\color{Accent}}{}{0pt}{}
\titleformat{\subsubsection}{\normalsize\bfseries\color{Ink}}{}{0pt}{}
\titlespacing*{\section}{0pt}{2.2ex plus 1ex minus .2ex}{1ex}
\titlespacing*{\subsection}{0pt}{1.8ex plus .6ex}{.7ex}
\pagestyle{fancy}
\fancyhf{}
\fancyhead[L]{\small\sffamily\color{Muted} AIYAGARI / FORMAL THEORY}
\fancyhead[R]{\small\sffamily\color{Muted} ARCHITECTURE v1.0}
\fancyfoot[L]{\footnotesize\color{Muted} Design specification — not a completed Lean formalization}
\fancyfoot[R]{\small\thepage}
\renewcommand{\headrulewidth}{0.3pt}
\renewcommand{\footrulewidth}{0pt}
\setlength{\headheight}{15pt}
\setlength{\emergencystretch}{3em}
\setlength{\parindent}{0pt}
\setlength{\parskip}{5.5pt plus 1pt}
\setlength{\LTpre}{8pt}
\setlength{\LTpost}{8pt}
\RecustomVerbatimEnvironment{verbatim}{Verbatim}{breaklines=true,fontsize=\small,breakanywhere=true}
\AtBeginEnvironment{longtable}{\small}
\widowpenalty=10000
\clubpenalty=10000
\ifPDFTeX
\DeclareUnicodeCharacter{2192}{\ensuremath{\to}}
\DeclareUnicodeCharacter{21A6}{\ensuremath{\mapsto}}
\DeclareUnicodeCharacter{00D7}{\ensuremath{\times}}
\fi
\ifLuaTeX
  \usepackage{selnolig}  % disable illegal ligatures
\fi
\usepackage{bookmark}
\IfFileExists{xurl.sty}{\usepackage{xurl}}{} % add URL line breaks if available
\urlstyle{same}
\hypersetup{
  pdftitle={Aiyagari (1993/1994)},
  colorlinks=true,
  linkcolor={Accent},
  filecolor={Maroon},
  citecolor={Accent},
  urlcolor={Accent},
  pdfcreator={LaTeX via pandoc}}

\title{Aiyagari (1993/1994)}
\usepackage{etoolbox}
\makeatletter
\providecommand{\subtitle}[1]{% add subtitle to \maketitle
  \apptocmd{\@title}{\par {\large #1 \par}}{}{}
}
\makeatother
\subtitle{Theory-only Lean replication architecture and Codex execution
specification}
\author{}
\date{September 11, 2026 -\/- Version 1.0}

\begin{document}
\maketitle

{
\hypersetup{linkcolor=}
\setcounter{tocdepth}{2}
\tableofcontents
}
\section{Status and use}\label{status-and-use}

This is a mathematical specification and an execution package, not a
completed Lean formalization. All theorem targets begin
\textbf{UNFORMALIZED}. The arguments below are proposed proof
constructions; they have not been checked by Lean. A successful build
alone will not establish economic adequacy. The executable prompts
implement this specification in small, reviewed milestones.

The published target is Aiyagari (1994), with the analytic model and
appendix in the December 1993 Working Paper 502. The primary release is
the continuous-asset, i.i.d.-income, bounded-utility theory: household
optimization, Propositions 1--5 of the working paper with necessary
qualifications, stationary asset supply, stationary general-equilibrium
existence, and the interest-rate/capital/gross-saving comparison. Exact
debt and borrowing-limit identities are a separate extension layer.
There is no numerical solver, asset grid, empirical calibration,
numerical certificate, or Python economic model. Python files in this
package only extract sources and validate project metadata.

The bounded-utility release is not a proof for unbounded log/CRRA
utility or serially correlated earnings. Those require distinct theorem
contracts. Likewise, nonexistence of an invariant law, divergence of
stationary means as a parameter approaches a boundary, and almost-sure
divergence along a household path are different claims. No one of these
is substituted for another.

\section{1. Source audit and changes to the earlier
plan}\label{source-audit-and-changes-to-the-earlier-plan}

The source inventory records original archive members, SHA-256 hashes,
printed pages and PDF pages. The supplied file
\texttt{schechtman\_escudero\_1976.pdf} is the \textbf{1977} published
article. The original Aiyagari archive is used only for its four
source-paper PDFs; none of its Lean code, theory summaries, proof
documents, or audits is used. The Mortensen--Pissarides archive supplies
formatting conventions only.

\subsection{1.1 Exact external
dependencies}\label{exact-external-dependencies}

\begin{longtable}[]{@{}
  >{\raggedright\arraybackslash}p{(\columnwidth - 4\tabcolsep) * \real{0.3333}}
  >{\raggedright\arraybackslash}p{(\columnwidth - 4\tabcolsep) * \real{0.3333}}
  >{\raggedright\arraybackslash}p{(\columnwidth - 4\tabcolsep) * \real{0.3333}}@{}}
\toprule\noalign{}
\begin{minipage}[b]{\linewidth}\raggedright
Source
\end{minipage} & \begin{minipage}[b]{\linewidth}\raggedright
Inspected location
\end{minipage} & \begin{minipage}[b]{\linewidth}\raggedright
Use and limitation
\end{minipage} \\
\midrule\noalign{}
\endhead
\bottomrule\noalign{}
\endlastfoot
Aiyagari (1993) & Printed pp.~11--22 and 37--40; PDF pp.~12--23 and
38--41 & Model, economic conclusions, Appendix Propositions 1--5. The
PDF is scanned; mathematical passages were inspected as page images. \\
Aiyagari (1994) & Printed pp.~665--674; PDF pp.~8--17 & Published
theoretical target, especially equations (1)--(8) and notes 16,
18--29. \\
Schechtman--Escudero (1977) & Theorems 3.8--3.9, printed pp.~161--162;
PDF pp.~11--12 & Marginal-value ratio argument for upper drift. Aiyagari
uses an eventual bound on relative risk aversion rather than requiring
an asymptotic exponent. \\
Clarida (1987) & Section 2 and Appendix; printed pp.~340--344, 348--350
& Income-fluctuation problem with borrowing. Its distributional and
utility restrictions must not be transferred silently. \\
Clarida (1990) & Proposition 2.4, printed p.~548; PDF p.~7 & States
continuity and both boundary limits \textbf{without proof}, referring to
Bewley (1984). It is a source locator, not a proved Lean dependency. \\
Benveniste--Scheinkman (1979) & Lemma 1, printed p.~728; PDF p.~3 &
Concave-function differentiability from a differentiable lower touching
function. Formalize the one-dimensional lemma directly; do not assume an
envelope equation. \\
Stokey--Lucas--Prescott (1989), SLP & Chapter 12, especially Assumption
12.1, Lemma 12.11, Theorems 12.12--12.13, printed pp.~381--385; PDF
pp.~391--395 & Endpoint crossing, oscillation contraction,
uniqueness/weak stability, parameter continuity. Definitions and proofs
are available in the supplied full book. \\
Chamberlain--Wilson (2000) & Sections 2--4, especially Theorem
4/Corollary 2; also Theorem 7 & Later published version of a dependency
cited by Aiyagari as a 1984 working paper. Guides the value-marginal
approach and the separately scoped pathwise result. \\
Sibley (1975); Miller (1976) & Supplied complete articles & Risk-order
claims are separately reviewed. Do not infer an aggregate ordering from
an ordering of household policies. \\
\end{longtable}

No additional paper is a prerequisite for the core execution sequence.
In particular, the missing Bewley manuscript is bypassed by the explicit
nonstationarity/tightness construction in Sections 9--10. This
construction must be formalized, not treated as a replacement citation.

\subsection{1.2 Corrections that must survive
implementation}\label{corrections-that-must-survive-implementation}

\begin{enumerate}
\def\labelenumi{\arabic{enumi}.}
\tightlist
\item
  \textbf{No-Ponzi is not zero terminal wealth.} Aiyagari (1993),
  Proposition 1, uses a nonnegative discounted-wealth limit. A feasible,
  nonoptimal household can leave positive discounted terminal wealth.
  Proving that the limit equals zero is a different optimality
  statement.
\item
  \textbf{A Bellman fixed point is not yet a lifetime optimum.}
  Verification against all nonanticipative feasible plans is an explicit
  theorem. Verification only against stationary policies is
  insufficient.
\item
  \textbf{Bounded utility does not bound the optimal consumption
  function.} The state space is unbounded. The drift argument actually
  requires consumption to become arbitrarily large as resources
  increase.
\item
  \textbf{Inada at zero plus a zero minimum effective income does not
  alone imply that the borrowing constraint never binds.} The note after
  Aiyagari (1993), Proposition 3, needs a further condition. An atom at
  zero combined with Inada is a sufficient condition; a continuous
  uniform effective-income distribution supplies a counterexample to the
  unqualified note. Section 5.3 gives an exact witness.
\item
  \textbf{Mixing is proved, not read off a diagram.} The lower
  transition map decreases strictly above the minimum; finite strings of
  near-minimum shocks and a high final shock provide SLP's
  common-horizon crossing probabilities.
\item
  \textbf{Local parameter continuity is not global operator-norm
  continuity on an unbounded state space.} Use finite-horizon continuity
  and a uniform discounted tail bound. Derive common compact invariant
  bounds only away from the critical interest rate.
\item
  \textbf{Stationary means do not follow from weak convergence alone.}
  Use a common compact state interval for interior parameter continuity.
  At the critical boundary, use bounded means only to obtain tightness,
  not to pass an unbounded integral through a weak limit.
\item
  \textbf{Never define equilibrium to require \(\beta(1+r)<1\).}
  Otherwise the headline interest-rate theorem becomes tautological.
  Construct the household optimum for every admissible positive gross
  return before imposing impatience for stationary-law existence.
\item
  \textbf{Fixed-limit equilibrium existence need not assume a bracketing
  condition.} Section 11 derives a lower bracket from the stationary
  household budget and the firm's production function.
\item
  \textbf{Net saving is zero in a stationary economy without growth.}
  The paper's higher saving-rate conclusion concerns gross investment
  \(\delta K/f(K)\).
\end{enumerate}

\section{2. Primitive model and normalized
notation}\label{primitive-model-and-normalized-notation}

At the base-model level, allow a probability measure \(\nu\) on a
compact labor interval \(L=[\ell_-,\ell_+]\) with
\(0<\ell_-\leq\ell_+<\infty\). The risky core adds \(\ell_-<\ell_+\) and
essential endpoints: every relative neighborhood of either endpoint has
positive probability. The deterministic benchmark is a separately
instantiated degenerate income law. Atoms are allowed; a density is not
required. In general-equilibrium modules impose \(\int\ell\,d\nu=1\).
Future labor draws are i.i.d. and independent of predetermined assets. A
household chooses after observing current labor income.

Let \(0<\beta<1\), \(R=1+r>0\), \(w>0\), and \(\phi\geq0\) with
\(w\ell_- -r\phi\geq0\). Set \[
e(\ell)=w\ell-r\phi,\qquad e_-=w\ell_- -r\phi,\qquad e_+=w\ell_+-r\phi.
\] The household's unshifted budget and normalized budget are \[
c_t+a_{t+1}=R a_t+w\ell_t,\quad a_{t+1}\geq-\phi,
\] \[
\widehat a_t=a_t+\phi,\quad z_t=R\widehat a_t+w\ell_t-r\phi,
\quad c_t+\widehat a_{t+1}=z_t,
\quad z_{t+1}=R\widehat a_{t+1}+e(\ell_{t+1}).
\] Use the fixed analytic state space \(X=\mathbb R_+\), not a
parameter-dependent interval. Economically reachable states satisfy
\(z\geq e_-\). The extension to \(0\leq z<e_-\) is a well-defined
auxiliary household problem and does not alter reachable-state behavior.

\begin{longtable}[]{@{}
  >{\raggedright\arraybackslash}p{(\columnwidth - 4\tabcolsep) * \real{0.3333}}
  >{\raggedright\arraybackslash}p{(\columnwidth - 4\tabcolsep) * \real{0.3333}}
  >{\raggedright\arraybackslash}p{(\columnwidth - 4\tabcolsep) * \real{0.3333}}@{}}
\toprule\noalign{}
\begin{minipage}[b]{\linewidth}\raggedright
Paper notation
\end{minipage} & \begin{minipage}[b]{\linewidth}\raggedright
Canonical mathematical meaning
\end{minipage} & \begin{minipage}[b]{\linewidth}\raggedright
Lean-facing name
\end{minipage} \\
\midrule\noalign{}
\endhead
\bottomrule\noalign{}
\endlastfoot
\(a_t\) & Beginning-of-period net asset holdings & \texttt{netAssets} \\
\(\widehat a_t\) & Shifted nonnegative holdings &
\texttt{shiftedAssets} \\
\(A(z)\) & Next-period \textbf{shifted}, not net, assets &
\texttt{assetPolicy} \\
\(c(z)=z-A(z)\) & Optimal consumption & \texttt{consumptionPolicy} \\
\(R=1+r\) & Gross return & \texttt{grossReturn} \\
\(\lambda=\beta^{-1}-1\) & Rate of time preference &
\texttt{impatienceRate} \\
\(\phi\) & Effective debt limit & \texttt{effectiveLimit} \\
\(V\) & Canonical bounded Bellman fixed point &
\texttt{valueFunction} \\
\(q(z)=V'_+(z)\) & Right marginal value; not automatically \(U'(0)\) &
\texttt{rightMarginalValue} \\
\(P\) & Kernel \(z\mapsto(RA(z)+e)_\#\nu\) & \texttt{householdKernel} \\
\(\pi\) & Probability law of total resources & \texttt{stationaryLaw} \\
\(\rho\) & Law of net assets, \((A-\phi)_\#\pi\) &
\texttt{netAssetLaw} \\
\(S(r,w,\phi)\) & \(\int(A(z)-\phi)\,d\pi(z)\) &
\texttt{stationaryAssetSupply} \\
\end{longtable}

\subsection{Utility assumptions}\label{utility-assumptions}

The core utility is real-valued and continuous on \(\mathbb R_+\),
bounded above and below, strictly increasing, strictly concave, and
twice continuously differentiable on \((0,\infty)\) with \(U'>0\) and
\(U''\leq0\). Differentiability is \textbf{not} required at zero. The
right marginal at zero may be finite or infinite. For drift only,
require constants \(C_0>0\) and \(M<\infty\) such that \[
-cU''(c)/U'(c)\leq M\quad(c\geq C_0).
\] The \(C^2\) condition is an explicit analytic regularity choice for
the curvature-based branch, not a hidden consequence of concavity.
Modules needing only continuity/concavity must not import it
unnecessarily. A nonempty primitive example is \(U(c)=c/(1+c)\), with
relative risk aversion \(2c/(1+c)\leq2\). The Inada diagnostic uses
\(U(c)=\sqrt c/(1+\sqrt c)\).

\subsection{Two borrowing
specifications}\label{two-borrowing-specifications}

For a finite institutional limit \(b\geq0\): \[
\phi_b(w,r)=\begin{cases}\min\{b,w\ell_-/r\},&r>0,\\b,&-1<r\leq0.\end{cases}
\] For the present-value limit, use \(r>0\) and
\(\phi_N(w,r)=w\ell_-/r\). These are separate model families. Continuity
of \(\phi_b\) at \(r=0\) must be proved using the locally binding finite
cap, not by evaluating \(1/r\) at zero.

\subsection{Assumption discipline}\label{assumption-discipline}

Primitive records may contain only economic data, probability
normalization/support, utility regularity, and production regularity.
They may not contain a value function, a savings rule, compact absorbing
bounds, mixing, an invariant law, asset-supply continuity, or
equilibrium existence. A general mathematical lemma can legitimately
assume that an arbitrary kernel is Feller and has a crossing property;
the economic wrapper must supply \textbf{proved} instances of these
hypotheses. Similarly, optimality and invariance are legitimate defining
conditions of an equilibrium object, but an existence theorem must
construct an object satisfying them.

\section{3. Bellman construction and lifetime
verification}\label{bellman-construction-and-lifetime-verification}

For \(v\in C_b(X)\) define \[
(Tv)(z)=\max_{0\leq a\leq z}\left\{U(z-a)+\beta\int v(Ra+e(\ell))\,d\nu(\ell)\right\}.
\] For continuity arguments only, reparameterize \(a=tz\),
\(t\in[0,1]\). The feasible action itself remains \(a\in[0,z]\); \(t\)
is not a uniquely selected control at \(z=0\).

\textbf{Existence and contraction.} The objective is continuous on each
compact action interval. Boundedness of \(U\) and \(v\) bounds \(Tv\).
On a compact neighborhood of a state, the transition image of actions
and labor states is compact, giving continuity of the integral and of
the maximum. Prove \[
\|Tv-T\widetilde v\|_\infty\leq\beta\|v-\widetilde v\|_\infty.
\] Use Mathlib's Banach fixed-point theorem to define \(V\). Its
uniqueness concerns bounded continuous fixed points. Concavity follows
by preservation under \(T\) and uniform convergence of iterations
starting from zero. Increasingness follows by the same preservation
argument; strict increasingness follows by retaining the old savings
action and consuming the additional resources.

\textbf{Policy.} The objective is strictly concave in \(a\) because
\(U(z-a)\) is strictly concave and the continuation value is concave.
Its unique maximizer defines \(A(z)\). Compact maximization and
uniqueness give continuity; at zero use \(0\leq A(z)\leq z\). The policy
is not an input to \texttt{HouseholdPrimitives}.

\textbf{Verification.} A feasible plan is a sequence of measurable
choices depending on the history observed so far, satisfying the shifted
budget and nonnegativity. Prove the finite-horizon Bellman inequality \[
V(z_0)\geq\mathbb E\!\left[\sum_{t=0}^{n-1}\beta^tU(c_t)+\beta^nV(z_n)\right]
\] for every such plan, with equality for the canonical policy. This can
be done using finite product measures of labor histories; an infinite
sample-path construction is not needed for this theorem. Since \(V\) is
bounded, the terminal term vanishes uniformly. The bounded discounted
utility series converges absolutely in expectation. Conclude that the
canonical policy solves the original infinite-horizon objective over all
nonanticipative feasible plans. Include the budget-change-of-variables
equivalence, rather than verifying only a newly invented problem.

\section{4. Policy shape, right marginals, and the envelope
theorem}\label{policy-shape-right-marginals-and-the-envelope-theorem}

Write \(G(a)=\beta\int V(Ra+e)\,d\nu\). It is concave. Comparing the
optimality inequalities for two resource levels, and using concavity of
\(U\) and \(G\), proves \[
0\leq A(z_2)-A(z_1)\leq z_2-z_1,\qquad
0\leq c(z_2)-c(z_1)\leq z_2-z_1\quad(z_1\leq z_2).
\] Do not replace these order statements with a differentiability claim
about the policy. Strict order claims in the source are tracked
separately; weak order and the Lipschitz bounds suffice for the core
proof chain.

\subsection{4.1 A marginal inequality valid before imposing
impatience}\label{a-marginal-inequality-valid-before-imposing-impatience}

For every \(z>0\), a finite, strictly positive right derivative
\(q(z)=V'_+(z)\) exists by concavity and strict increasingness. It is
nonincreasing. If \(D=(\sup U-\inf U)/(1-\beta)\), then \[
0<q(z)\leq\frac{V(z)-V(0)}{z}\leq D/z.
\] Retain the optimal current consumption and invest an additional
initial amount \(h>0\). This gives \[
V(z+h)-V(z)\geq\beta\int[V(RA(z)+e+Rh)-V(RA(z)+e)]\,d\nu.
\] Divide by \(h\) and let \(h\downarrow0\). Nonnegative concave
difference quotients permit monotone convergence, yielding \[
q(z)\geq\beta R\int q(RA(z)+e)\,d\nu. \tag{M}
\] Boundary derivatives at zero are treated as extended limits until
their finiteness or null relevance has been proved. Do not use Lean's
real-valued integral to turn an infinite marginal expectation into zero.
Inequality (M) itself establishes the required conditional integrability
at states with finite \(q\).

\subsection{4.2 Consumption positivity in the impatient
region}\label{consumption-positivity-in-the-impatient-region}

Under \(\beta R<1\), prove \(c(z)>0\) for \(z>0\).

If \(U'_+(0)=\infty\), consuming an increment of additional resources
while keeping the original savings choice would imply \(q(z)=\infty\)
whenever \(c(z)=0\), contradicting finite \(q(z)\).

If \(U'_+(0)=L<\infty\), prove first, by finite-horizon induction, that
\(V\) is \(L\)-Lipschitz when \(\beta R\leq1\). For a larger state
\(z+h\) with an optimal action \(a\), compare it with action \(a\) at
\(z\) when \(a\leq z\), and with action \(z\) when \(a>z\). In the
latter case the extra resources split between extra current consumption
and extra savings; the return-adjusted continuation slope is at most
\(\beta RL\leq L\). Passing to the limit proves the bound. If
\(c(z)=0\), transferring a small amount from savings to consumption
gains a marginal amount approaching \(L\), while the continuation loss
is at most \(\beta RL<L\). This contradicts optimality.

\subsection{4.3 Envelope and Euler
conditions}\label{envelope-and-euler-conditions}

At any \(z>0\) with \(c(z)>0\), keep \(A(z)\) fixed while perturbing
\(z\) in a small two-sided neighborhood. The differentiable concave
function \[
W(x)=U(x-A(z))+\beta\int V(RA(z)+e)\,d\nu
\] touches \(V\) from below. The one-dimensional Benveniste--Scheinkman
lemma gives \[
V'(z)=U'(c(z)).
\] Prove the touching lemma from one-sided slopes; importing a
multidimensional subgradient library is unnecessary. Continuity of \(c\)
and \(U'\) then gives continuity of \(V'\) on the positive-consumption
region.

At \(A(z)>0\), nearby savings choices keep next resources bounded away
from zero, permitting differentiation under the integral using a local
marginal bound. First-order optimality yields equality in (M). At
\(A(z)=0\), use one-sided difference quotients; they yield the Euler
inequality and rule out an infinite continuation marginal when current
consumption is positive. State the Euler result with every integrability
obligation visible.

\section{5. Borrowing thresholds and an important
diagnostic}\label{borrowing-thresholds-and-an-important-diagnostic}

\subsection{5.1 A positive binding
interval}\label{a-positive-binding-interval}

If \(\beta R<1\) and either \(e_->0\) or \(U'_+(0)<\infty\), then
\(q(e_-)\) is finite and positive. If \(A(z)>0\), the Euler equality and
monotonicity give \[
q(z)=\beta R\mathbb E q(RA(z)+e)\leq\beta Rq(e_-).
\] Right continuity of the right derivative at \(e_-\) makes this
impossible for all \(z\) in a sufficiently small right neighborhood of
\(e_-\). Consequently, some \(\widehat z>e_-\) satisfies \(A(z)=0\) on
\([e_-,\widehat z]\). This reproduces the qualified Proposition 3. Do
not assume the threshold exists before this argument.

\subsection{5.2 A sufficient condition for never
binding}\label{a-sufficient-condition-for-never-binding}

If \(e_-=0\), \(U'_+(0)=\infty\), and the income distribution has an
atom at \(e=0\), then \(A(z)>0\) for every \(z>0\). A small positive
saving deviation from \(A=0\) gains at least the zero-income probability
times the increase in \(V\) near zero, whose right slope is infinite;
its current-utility cost has finite slope. More generally, an infinite
expected right marginal of continuation value at \(A=0\) is sufficient.
Inada and \(e_-=0\) by themselves are not sufficient.

\subsection{5.3 Exact continuous-state counterexample to the unqualified
note}\label{exact-continuous-state-counterexample-to-the-unqualified-note}

Take \(\beta=1/2\), \(R=3/2\), and effective income uniform on
\([0,1]\). Use \[
U(c)=\frac{\sqrt c}{1+\sqrt c},\qquad U(0)=0.
\] This utility is bounded, strictly increasing, strictly concave, Inada
at zero, and has relative risk aversion \[
-\frac{cU''(c)}{U'(c)}=\frac12+\frac{\sqrt c}{1+\sqrt c}<\frac32.
\] The setup is economically admissible in the original model: labor
uniform on \([2/3,4/3]\) has mean one; choose \(w=3/2\), \(r=1/2\) and
natural limit \(\phi=2\).

Because \(V\) is bounded and continuous, \[
\left.\frac{d}{da^+}\int_0^1V(Ra+e)\,de\right|_{a=0}
=R[V(1)-V(0)].
\] Also \(0\leq V\leq2\). Thus the right derivative of the Bellman
objective at \(a=0\) is at most \(-U'(z)+3/2\). For \(0<z\leq1/100\),
\(U'(z)\geq U'(1/100)>3/2\). Concavity therefore makes \(A(z)=0\)
optimal. This disproves the unqualified nonbinding note without using a
finite asset grid. The counterexample must itself be verified in Lean
before it is labeled a certified correction.

\section{6. Joint parameter continuity and uniform upper
drift}\label{joint-parameter-continuity-and-uniform-upper-drift}

\subsection{6.1 Continuity before
stationarity}\label{continuity-before-stationarity}

Keep \(U\) and \(\nu\) fixed. First use normalized prices
\(\theta=(R,w,k)\) with \(e(\ell)=w\ell+k\geq0\); the original
coordinates \((w,R,\phi)\) map to these by \(k=-(R-1)\phi\). Let
\(\theta\) vary over this admissible normalized domain. Each
finite-horizon value function \(T_\theta^n0\) is jointly continuous in
\((\theta,z)\), by the fixed \(t\in[0,1]\) maximization representation
and compactness of the shock space. Its distance from \(V_\theta\) is
bounded uniformly in \(\theta\) by a constant times \(\beta^n\).
Consequently, \(V_\theta(z)\) is jointly continuous. The
unique-maximizer argument gives joint continuity of \(A_\theta(z)\).
This includes parameter values with \(\beta R=1\) and \(\beta R>1\): the
bounded-utility Bellman construction does not require impatience.

\subsection{6.2 A locally uniform version of Proposition
4}\label{a-locally-uniform-version-of-proposition-4}

Work on a parameter neighborhood with \[
0<R_{\min}\leq R\leq R_{\max},\quad \beta R\leq\gamma_*<1,
\quad 0\leq e_-\leq e_+\leq E_*,\quad e_+-e_-\leq\Delta_*.
\] Replace the eventual relative-risk-aversion bound by a larger
positive integer \(m\). For \(c_2\geq c_1\geq C_0\), differentiation of
\(c^mU'(c)\) proves \[
\frac{U'(c_1)}{U'(c_2)}\leq(c_2/c_1)^m.
\] Choose \(C\geq C_0\) sufficiently large that \[
\gamma_*\left(1+\Delta_*/C\right)^m<1.
\] The bound \(q_\theta(x)\leq D/x\) and the envelope identity imply
that \(c_\theta(x)>C\) whenever \(x\) exceeds a common level
\(L>D/U'(C)\). Choose \(K\) with \(R_{\min}K>L\), and then
\(B\geq R_{\max}K+E_*\).

If \(A_\theta(z)\leq K\), then \(RA_\theta(z)+e_+\leq B\leq z\) for
\(z\geq B\). Otherwise both extreme next-period states have consumption
exceeding \(C\). The consumption Lipschitz bound and curvature
inequality imply \[
\frac{\mathbb E q_\theta(RA_\theta(z)+e)}{q_\theta(RA_\theta(z)+e_+)}
\leq(1+\Delta_*/C)^m.
\] Euler equality then gives \(q_\theta(z)<q_\theta(RA_\theta(z)+e_+)\).
Since \(q_\theta\) is nonincreasing, \[
RA_\theta(z)+e_+<z.
\] This proves a uniform upper drift bound and a common
forward-invariant interval. It avoids an unsupported assertion that
pointwise absorbing bounds vary continuously in prices.

\textbf{Do not infer finite-time entry into \([e_-,B]\) from weak
downward drift.} A trajectory can approach a boundary asymptotically.
Global weak convergence will be obtained by applying the same
compact-state theorem on larger intervals, not by assuming finite-time
entry.

\section{7. Kernel, derived crossing, and the SLP
theorem}\label{kernel-derived-crossing-and-the-slp-theorem}

Define the Markov kernel by the pushforward of \(\nu\) under
\(\ell\mapsto RA(z)+e(\ell)\). Verify measurability, probability mass
one, and the integral identity for bounded measurable test functions.
Joint continuity of the transition map implies the Feller property.
Monotonicity of \(A\) proves stochastic monotonicity by using the same
labor realization at both initial states.

\subsection{7.1 Deriving mixing from household
optimality}\label{deriving-mixing-from-household-optimality}

On a compact invariant interval \(I=[e_-,B]\), define \[
h_-(z)=RA(z)+e_-.
\] For \(z>e_-\), if \(h_-(z)\geq z\), then \(A(z)>0\) and Euler
equality gives \[
q(z)=\beta R\mathbb E q(RA(z)+e)\leq\beta Rq(h_-(z))
\leq\beta Rq(z)<q(z),
\] a contradiction. Thus \(h_-(z)<z\) above \(e_-\). At the lower
endpoint, \(A(e_-)=0\): this is immediate if \(e_-=0\), and otherwise
follows from the same inequality. Repeated application of \(h_-\) to
\(B\) decreases to \(e_-\) by continuity.

Choose \(d\in(e_-,e_+)\). There is a finite \(N\geq1\) with
\(h_-^N(B)<d\). By continuity of this finite composition, there is an
\(\eta>0\) such that \(N\) shocks all in \([e_-,e_-+\eta]\) also move
\(B\) below \(d\). Their probability is \(p_-^N>0\). Regardless of the
earlier history, a final shock above \(d\) puts the state above \(d\),
with probability \(p_+>0\). Choose
\(\varepsilon=\min\{p_-^N,p_+,1/2\}\). Then \[
P^N(e_-,[d,B])\geq\varepsilon,
\qquad P^N(B,[e_-,d])\geq\varepsilon.
\] This is precisely the needed common-horizon crossing condition. It
does not require an atom at the lower endpoint, a positive density, a
unique upper fixed point, or a drawing supplied by the paper.

\subsection{7.2 Formalizing the one-dimensional SLP
argument}\label{formalizing-the-one-dimensional-slp-argument}

First prove compact monotone-Feller invariant existence using the
increasing measures \(P^{*n}\delta_{e_-}\) and decreasing measures
\(P^{*n}\delta_B\). Compactness of probability measures supplies cluster
points; order-determining continuous increasing tests give convergence;
Feller continuity makes both limits invariant.

For a bounded nondecreasing test function \(f\), crossing gives \[
\varepsilon f(d)+(1-\varepsilon)f(e_-)
\leq P^Nf(z)
\leq\varepsilon f(d)+(1-\varepsilon)f(B).
\] Apply this to \(P^{kN}f\) and induct: \[
(P^{kN}f)(B)-(P^{kN}f)(e_-)
\leq(1-\varepsilon)^k[f(B)-f(e_-)].
\] The two endpoint invariant limits therefore agree on continuous
increasing tests and hence as probability measures. Every initial law
lies between the endpoint laws in stochastic order, proving uniqueness
and \textbf{weak} convergence. No total-variation convergence is
claimed.

For a point \(z>B\), run the compact theorem on \([e_-,\max\{B,z\}]\);
the previously constructed invariant law also lives there, so uniqueness
identifies the limit. Initial states below \(e_-\) enter the economic
state interval after one step. Dominated convergence for bounded tests
then extends convergence to every initial probability law on
\(\mathbb R_+\). This also proves uniqueness of an invariant law on the
full state space, including laws not initially known to have finite
moments.

\section{8. Stationary continuity, aggregation, and the cross-sectional
bridge}\label{stationary-continuity-aggregation-and-the-cross-sectional-bridge}

Use the uniform drift result locally around a strictly impatient price
vector. Regard all invariant laws as measures on the same compact
interval \([0,B]\) (the fixed lower endpoint zero avoids
varying-state-space problems). Joint kernel continuity and uniqueness
imply weak continuity of the stationary law, by the subsequence argument
of SLP Theorem 12.13. Show explicitly that every subsequential limit is
invariant under the limiting kernel.

For average assets, split \[
\left|\int A_{\theta_n}\,d\pi_{\theta_n}-\int A_{\theta_0}\,d\pi_{\theta_0}\right|
\leq\|A_{\theta_n}-A_{\theta_0}\|_{[0,B]}
+\left|\int A_{\theta_0}\,d\pi_{\theta_n}-\int A_{\theta_0}\,d\pi_{\theta_0}\right|.
\] Joint continuity gives uniform convergence on \([0,B]\); weak
convergence handles the fixed continuous integrand. Subtract the
continuous effective debt limit. This proves stationary asset-supply
continuity without an unproved moment assumption.

The canonical resource law and the economically meaningful asset/labor
cross-section must be linked. Let \(\rho=(A-\phi)_\#\pi\). Predetermined
assets and current i.i.d. labor have joint law \(\rho\otimes\nu\). Its
resource image is \(\pi\) by invariance. Prove this identity and then \[
\mathbb E_\pi z=R\mathbb E_\pi A+\mathbb E e,
\quad S=\mathbb E_\pi A-\phi,
\quad \mathbb E_\pi c=rS+w\mathbb E\ell.
\] In general equilibrium \(\mathbb E\ell=1\). These statements require
explicit integrability, supplied here by compact support. A continuum of
individually independent random variables is not needed: the equilibrium
is defined by a stationary cross-sectional law, not by an unproved
continuum law of large numbers.

\section{\texorpdfstring{9. Excluding a stationary law when
\(\beta R\geq1\)}{9. Excluding a stationary law when \textbackslash beta R\textbackslash geq1}}\label{excluding-a-stationary-law-when-beta-rgeq1}

This is a proposed replacement proof, inspired by the marginal-value
argument in Chamberlain--Wilson, rather than a claim that Clarida
supplies the missing derivation. The use of the right marginal of
\textbf{value} is important: \(U'(c)\) need not describe the value
marginal at a zero-consumption corner.

\subsection{9.1 Resolve the zero-resource state
first}\label{resolve-the-zero-resource-state-first}

The right derivative at zero may be finite or infinite. If it is finite,
it is positive and may be included in the real-valued function \(q\). If
it is infinite and \(p_0=\Pr(e=0)>0\), optimal savings from every
positive state must be strictly positive: a deviation from zero savings
would have infinite marginal continuation gain. Thus transitions from
positive states cannot hit zero. An invariant law then satisfies
\(\pi\{0\}=p_0\pi\{0\}\); nondegeneracy gives \(p_0<1\), hence
\(\pi\{0\}=0\). If \(p_0=0\), zero has no incoming probability. Only
after this proof may \(q\) be assigned an arbitrary finite value at an
irrelevant zero state.

\subsection{9.2 A bounded strict-concavity argument avoids stationary
marginal
integrability}\label{a-bounded-strict-concavity-argument-avoids-stationary-marginal-integrability}

Suppose \(\pi\) is invariant. Set \(\gamma=\beta R\geq1\) and \[
\psi(x)=\frac{x}{1+x},\qquad x\geq0.
\] For \(\pi\)-almost every state, (M) gives a finite conditional mean
\(m(z)=\int q(z')P(z,dz')\) and \[
\psi(q(z))\geq\psi(\gamma m(z))\geq\psi(m(z))
\geq\int\psi(q(z'))P(z,dz').
\] The end terms have the same integral by stationarity, since \(\psi\)
is bounded. If \(\gamma>1\), strict monotonicity and \(m(z)>0\) make the
middle inequality strict almost everywhere, a contradiction.

At \(\gamma=1\), all inequalities are equalities almost everywhere.
Strict Jensen equality gives \(q(z')=q(z)\) under the stationary
one-step joint law. A convenient elementary proof of strict Jensen uses
\[
\psi(m)+\psi'(m)(x-m)-\psi(x)
=\frac{(x-m)^2}{(1+m)^2(1+x)}.
\] The nonnegative tangent gap has zero conditional integral exactly
when \(x=m\) almost surely. Only conditional integrability of \(q\) is
used; \(\int q\,d\pi<\infty\) is \textbf{not} assumed.

\subsection{9.3 Equal value marginals imply equal
consumption}\label{equal-value-marginals-imply-equal-consumption}

At positive consumption, the envelope proof gives \(q=U'(c)\). At zero
consumption, consuming an increment of extra resources while retaining
savings gives \(q\geq U'_+(0)\). Strict concavity makes \(U'\) injective
on \((0,\infty)\) and \(U'(c)<U'_+(0)\) for \(c>0\). Therefore equal
finite value marginals imply equal consumption, including the
finite-marginal zero-consumption corner. Along any finite stationary
history at \(\gamma=1\), consumption is thus equal to its initial value.

\subsection{9.4 Two independent future shock strings give the
contradiction}\label{two-independent-future-shock-strings-give-the-contradiction}

Here \(R=1/\beta>1\). Draw a common initial state \(Z_0\sim\pi\) and two
independent length-\(n\) future effective-income strings \((e_j)\) and
\((\widetilde e_j)\), independent of \(Z_0\). Evolve both by the optimal
policy. Marginal stationarity gives \(Z_n,\widetilde Z_n\sim\pi\).
Constant consumption along both paths gives the exact cancellation \[
R^{-n}(Z_n-\widetilde Z_n)
=\sum_{j=1}^nR^{-j}(e_j-\widetilde e_j)=D_n.
\] Every probability law on \(\mathbb R_+\) is tight, so the left side
tends to zero in probability; use the union bound on \(Z_n\) and
\(\widetilde Z_n\), not a moment assumption. The right side is bounded
uniformly by \((e_+-e_-)/(R-1)\). Hence \(\mathbb E D_n^2\to0\).

On the other hand, independence and nondegenerate bounded income imply
\[
\mathbb E D_n^2=2\operatorname{Var}(e)\sum_{j=1}^nR^{-2j}
\geq2\operatorname{Var}(e)R^{-2}>0\quad(n\geq1),
\] a contradiction. The products can be constructed separately for each
\(n\); this proof does not require building an infinite stationary path
space.

Conclude that the canonical household kernel has no invariant
probability law when \(\beta R\geq1\). The theorem must not assume
bounded assets, finite stationary marginal utility, or a stationary law
supported away from zero.

\section{10. The two stationary asset-supply boundary
limits}\label{the-two-stationary-asset-supply-boundary-limits}

\subsection{\texorpdfstring{10.1 Divergence as
\(r\uparrow\lambda\)}{10.1 Divergence as r\textbackslash uparrow\textbackslash lambda}}\label{divergence-as-ruparrowlambda}

Allow \(w_n\to w_*>0\), \(r_n\uparrow\lambda\), and effective limits
\(\phi_n\to\phi_*<\infty\). Let \(\pi_n\) be the subcritical invariant
laws. If \(S_n\) does not tend to \(+\infty\), choose a subsequence with
\(S_n\leq M\). Then \[
\mathbb E_{\pi_n}A_n=S_n+\phi_n
\] is bounded. Invariance implies
\(\mathbb E_{\pi_n}z=R_n\mathbb E A_n+\mathbb E e_n\), so the
nonnegative resource means are uniformly bounded. Markov's inequality
gives tightness, and Prokhorov gives a weakly convergent subsequence.

For any bounded continuous test \(f\), joint policy/transition
continuity gives \(P_nf\to P_*f\) uniformly on every compact resource
interval. The tight tail controls the complement. Passing through the
stationarity equation using \textbf{bounded} tests proves that the limit
law is invariant for the critical kernel. Section 9 rules this out.
Therefore \(S_n\to+\infty\).

This proves the parameter-boundary result directly. It does not infer it
from exploding supports, almost-sure paths, or convergence of unbounded
integrals. State the sequence theorem first; obtain the one-sided limit
with Lean's filter API afterward.

\subsection{\texorpdfstring{10.2 Natural debt limit as
\(r\downarrow0\)}{10.2 Natural debt limit as r\textbackslash downarrow0}}\label{natural-debt-limit-as-rdownarrow0}

For \(\phi=w\ell_-/r\), effective income is \(e=w(\ell-\ell_-)\). If
\(w\to w_0>0\), then \(R\to1\) remains strictly impatient, and the
uniform drift proof bounds \(A\) on one common compact set. Thus
\(\mathbb EA\) stays bounded while \(\phi\to+\infty\), giving
\(S\to-\infty\).

The divergence of \(\phi\) does not obstruct the argument: apply the
household continuity/drift modules to the normalized income law, not to
an unbounded raw parameter triple \((w,r,\phi)\).

\section{11. Firms, existence, and the main economic
result}\label{firms-existence-and-the-main-economic-result}

A concrete production witness is \(f(K)=\sqrt K\) with \(\delta=1/2\);
combined with the household witness in Section 2, it demonstrates a
nonempty full primitive class.

Let \(f:[0,\infty)\to[0,\infty)\) be continuous with \(f(0)=0\), twice
continuously differentiable on \((0,\infty)\), \(f'>0\), \(f''<0\),
\(f'(0+)=\infty\), and \(f'(\infty)=0\). Let \(0<\delta<1\). The
constant-returns technology is \(F(K,L)=Lf(K/L)\) for \(L>0\).

For \(r>-\delta\), construct the unique capital demand \(K(r)>0\) from
\(f'(K(r))=r+\delta\) and define \[
w(r)=f(K(r))-K(r)f'(K(r)).
\] Prove these functions are continuous, \(K\) is strictly decreasing,
and \(w>0\). They are derived from \(f\), not free curves. The labor
normalization is one.

\subsection{11.1 A noncircular equilibrium
definition}\label{a-noncircular-equilibrium-definition}

A stationary equilibrium contains prices and a cross-sectional law
satisfying firm optimization, household lifetime optimality, invariance,
finite first moments needed for aggregation, and capital-market
clearing. The household optimum is the canonical Bellman solution
constructed for all \(R>0\). Neither the definition nor its primitive
hypotheses require \(r<\lambda\), asset-supply monotonicity, or
uniqueness of equilibrium. Prove equivalence between a resource-law
formulation and the net-asset/current-labor formulation.

\subsection{\texorpdfstring{11.2 A derived lower bracket for every
finite
\(b\)}{11.2 A derived lower bracket for every finite b}}\label{a-derived-lower-bracket-for-every-finite-b}

The firm assumptions imply \(f(K)/K\to0\) as \(K\to\infty\); prove this
with a tangent bound and \(f'(K)\to0\). Choose a large \(K_L\) such that
\(f'(K_L)<\delta\) and \(f(K_L)<\delta K_L\). Put
\(r_L=f'(K_L)-\delta\in(-\delta,0)\) and \(w_L=f(K_L)-K_Lf'(K_L)>0\).

For the household stationary law at these prices, nonnegative
consumption and the stationary budget give \[
0\leq\mathbb Ec=w_L+r_LS(r_L),
\quad S(r_L)\leq\frac{w_L}{-r_L}<K_L.
\] The last inequality is exactly \(f(K_L)<\delta K_L\). Thus excess
asset supply is negative at an explicitly derived feasible price. Near
\(\lambda\), asset supply tends to infinity while \(K(r)\) remains
finite. Continuity and the intermediate value theorem give at least one
stationary equilibrium for every finite institutional limit \(b\geq0\),
with \(r\in(-\delta,\lambda)\).

This is stronger than the earlier plan's conditional bracketing theorem,
without imposing a conclusion as an assumption. Equilibrium need not be
unique and its interest rate need not be positive.

\subsection{11.3 Natural-limit
equilibrium}\label{natural-limit-equilibrium}

At \(r\downarrow0\), capital demand and wages tend to finite positive
values because \(\delta>0\). Section 10.2 gives negative infinite asset
supply; Section 10.1 gives positive infinite supply at the other
endpoint. Continuity gives an equilibrium with \(0<r<\lambda\).

\subsection{11.4 Certainty benchmark and capital/saving
comparisons}\label{certainty-benchmark-and-capitalsaving-comparisons}

First consider deterministic labor at its mean,
\(\bar\ell=\int\ell\,d\nu\), and a strictly impatient return. Its
constant effective income is \(\bar e=w\bar\ell-r\phi_C\), where
\(\phi_C\) uses the certainty income rather than the risky minimum. The
deterministic transition \(h(z)=RA(z)+\bar e\) satisfies
\(h(\bar e)=\bar e\) and \(h(z)<z\) above \(\bar e\) by Section 7.1,
without requiring nondegenerate income. Every finite initial state
converges to \(\bar e\); bounded-test dominated convergence gives the
unique invariant law \(\delta_{\bar e}\). Consequently, certainty
stationary net assets equal \(-\phi_C\).

For the corresponding risky economy,
\(S=\mathbb EA-\phi_R\geq-\phi_R\geq-\phi_C\), because shifted assets
are nonnegative and the risky income floor is no greater than mean
income. This gives weakly higher stationary assets at every subcritical
admissible rate. At fixed positive wages, the upper-boundary theorem
gives \(S>0\geq-\phi_C\) for all rates sufficiently close to \(\lambda\)
from below, yielding the paper's qualified strict precautionary-assets
comparison. It does not imply a general ordering across two risky income
distributions. These are contracts A04 and A05, implemented with the
benchmark module in milestone 09.

For the mean-income certainty economy, the stationary
representative-agent benchmark is \(r^{FI}=\lambda\) and
\(K^{FI}=K(\lambda)\). Verify the constant-consumption plan using the
concavity inequality for utility and the deterministic present-value
budget; an Euler equality alone is not a verification theorem. The
certainty borrowing limit is generally \(\min\{b,w/r\}\), not the
risky-income limit with \(\ell_-\). Since benchmark capital is positive,
it is feasible under either nonnegative institutional debt cap.

Any stationary equilibrium of the risky economy must have \(r<\lambda\)
by Section 9, independently of the construction used to prove existence.
Strictly decreasing capital demand gives \[
K>K^{FI}.
\] Moreover, \[
\frac{d}{dK}\left(\frac{\delta K}{f(K)}\right)
=\frac{\delta[f(K)-Kf'(K)]}{f(K)^2}>0.
\] Hence the gross saving/investment share is higher than in the
certainty benchmark. Goods-market clearing follows from the household
stationary budget and the firm's factor-payment identity:
\(\mathbb Ec+\delta K=f(K)\).

\section{12. No-Ponzi and exact extension
layer}\label{no-ponzi-and-exact-extension-layer}

\subsection{12.1 No-Ponzi equivalence}\label{no-ponzi-equivalence}

For \(r>0\), bounded labor income and the natural lower bound imply,
pathwise on the event of feasibility, \[
\sum_{t=0}^{T}R^{-t}c_t
=Ra_0+\sum_{t=0}^{T}R^{-t}w\ell_t-R^{-T}a_{T+1}.
\] The consumption sums are nondecreasing and bounded above. Thus the
discounted asset sequence has a finite nonnegative limit. It need not
have limit zero.

For the converse, assets \(a_t\) are measurable before \(\ell_t\) is
drawn. If the natural limit is violated with positive probability, first
extract a fixed deficit \(\eta>0\) on a positive-probability event. A
sufficiently long string of incomes within a fixed neighborhood of
\(w\ell_-\) makes the discounted debt deficit exceed all possible
discounted subsequent income, even if future income always equals its
upper bound and consumption is zero. Independence gives this finite
event positive conditional probability. The discounted wealth limit is
then negative on a positive-probability event, contradicting no-Ponzi.
The proof must retain the factor \(w\) in the size of the low-labor
neighborhood. This is a statement about adapted feasible plans, not an
equivalence for arbitrary pathwise arrays with foresight.

\subsection{12.2 Extensions with complete algebraic
contracts}\label{extensions-with-complete-algebraic-contracts}

At \(r=0\), effective income is \(w\ell\) and the shifted household
problem is independent of \(b\). Therefore \(S(b_2)-S(b_1)=-(b_2-b_1)\)
at fixed wages. If two finite caps both exceed the natural debt limit at
a fixed positive interest rate, their effective limits and hence their
entire household problems coincide.

In a pure-exchange economy with government debt \(d\) and lump-sum tax
\(rd\), use \(c_t+a_{t+1}=Ra_t+w\ell_t-rd\). Under the tax-adjusted
natural bound \(a_t\geq d-w\ell_-/r\), the change of variables
\(x_t=a_t-d\) gives exactly the debt-free household problem and
equilibrium clearing \(\mathbb Ex=0\). Represent the taxed model with a
general real asset floor, which may be positive; do not force its
original debt-limit parameter to be nonnegative. The transformed
debt-free problem satisfies the core nonnegative-shift convention. This
proves the debt-neutrality correspondence in that particular borrowing
regime; it does not prove neutrality under a fixed borrowing cap.

For the monetary reinterpretation, specify the real money stock, its
gross real return, and the tax-financing identity before proving a
budget correspondence. Do not use the capital-economy existence theorem
for a different clearing condition without a separate argument.

For isoelastic utility and gross trend growth \(G\), the Euler
normalization has the factor \(\beta R G^{-\sigma}\) and the transformed
discount factor is \(\widetilde\beta=\beta G^{1-\sigma}\). This is an
algebraic identity, not an extension of the bounded-utility stationarity
theorem to CRRA. A full growth theorem requires its own
unbounded-utility analysis.

\subsection{12.3 Claims intentionally not promoted to core
theorems}\label{claims-intentionally-not-promoted-to-core-theorems}

The qualitative core does not prove global monotonicity of asset supply
in the interest rate or borrowing limit; uniqueness of equilibrium; an
aggregate risk-order theorem; a multiple-equilibrium witness; or all
pathwise statements in the paper's footnotes. The supercritical pathwise
result \(\beta R>1\) has a short route from (M), Markov's inequality and
Borel--Cantelli. The critical pathwise result \(\beta R=1\) requires a
separate adaptation of Chamberlain--Wilson's uncertainty argument. It
remains a named extension, not a hidden premise of asset-supply
divergence.

The stronger strict-increase description of both household policies
above a threshold also remains a separate source claim until strictness
has been established. The core uses the proved weak monotonicity and
Lipschitz bounds. The release label must therefore say \textbf{core
stationary theory}, with the source-coverage table attached, rather than
claiming that every theoretical sentence or utility specification has
been formalized.

\section{13. Lean architecture and implementation
policy}\label{lean-architecture-and-implementation-policy}

Use namespace \texttt{Aiyagari1994} and a clean project directory,
separate from the failed implementation. Suggested module families are
\texttt{Primitives}, \texttt{Budget}, \texttt{Analysis},
\texttt{Household}, \texttt{Stationary}, \texttt{Aggregate},
\texttt{Firms}, \texttt{Equilibrium}, \texttt{Diagnostics}, and
\texttt{Extensions}. The theorem manifest supplies the
declaration/module contracts. Names are project targets, not assertions
that matching Mathlib declarations already exist.

The environment milestone pins a mutually compatible Lean version and
Mathlib commit, records them, and proves small API examples before
economic coding. Do not choose an unverified version number from this
document. Public Mathlib documentation currently exposes Banach fixed
points, probability measures with weak topology, Markov kernels,
compactness of probability measures, and Prokhorov compactness. The
installed pinned checkout, not the moving website, is the authority for
exact signatures.

Use probability measures and pushforwards rather than informal
distributions. Every real integral must come with the relevant
integrability fact before it is interpreted economically. Classical
choice is permitted only after a proved existence result. Generic
analytic theorems may take their mathematical hypotheses as arguments;
paper-level wrappers must discharge them from primitives.

No \texttt{sorry}, \texttt{admit}, custom \texttt{axiom}, kernel-check
bypass, or unapproved replacement theorem is allowed in completed
modules. Standard Lean foundations such as \texttt{propext},
\texttt{Classical.choice}, and \texttt{Quot.sound} are not defects;
audit the exact transitive axiom list instead of advertising
`axiom-free' proofs. \texttt{assert\_no\_sorry} catches
\texttt{sorryAx}, not economically circular assumptions. Manual adequacy
review remains necessary.

\texttt{All.lean} imports only completed substantive modules.
\texttt{Audit.lean} imports \texttt{All.lean} and checks every exported
contract declaration with \texttt{\#check}, \texttt{assert\_no\_sorry},
and \texttt{\#print\ axioms}. Unfinished targets remain in the manifest,
not in the build as placeholders. Staging a file in an `experimental'
directory is not a way to hide completed theorems' dependencies.

\section{14. Milestone sequence and review
gates}\label{milestone-sequence-and-review-gates}

\begin{longtable}[]{@{}
  >{\raggedright\arraybackslash}p{(\columnwidth - 4\tabcolsep) * \real{0.0886}}
  >{\raggedright\arraybackslash}p{(\columnwidth - 4\tabcolsep) * \real{0.6203}}
  >{\raggedright\arraybackslash}p{(\columnwidth - 4\tabcolsep) * \real{0.2911}}@{}}
\toprule\noalign{}
\begin{minipage}[b]{\linewidth}\raggedright
Prompt
\end{minipage} & \begin{minipage}[b]{\linewidth}\raggedright
Task
\end{minipage} & \begin{minipage}[b]{\linewidth}\raggedright
Recommended reasoning
\end{minipage} \\
\midrule\noalign{}
\endhead
\bottomrule\noalign{}
\endlastfoot
00 & Source hashes, clean environment, pinned toolchain, API probes &
High \\
01 & Primitive records, normalized budget, consistency witness & High \\
02 & Bellman construction and lifetime verification & Extra-high \\
03 & Policy order, right marginals, envelope/Euler, borrowing correction
& Extra-high \\
04 & Joint parameter continuity and uniform upper drift & Extra-high \\
05 & Kernel, derived mixing, compact SLP theorem, global weak stability
& Extra-high \\
06 & Invariant-law continuity, cross-sectional bridge, aggregation &
Extra-high \\
07a & Zero-state audit and bounded-Jensen marginal argument &
Extra-high \\
07b & Critical nonstationarity by the two-shock-string argument &
Extra-high \\
08 & Asset-supply boundary limits by tightness and normalized limits &
Extra-high \\
09 & Firms, both existence theorems, benchmark, headline comparisons &
Extra-high \\
10 & No-Ponzi equivalence and exact borrowing/debt extensions &
Extra-high for no-Ponzi; High for algebra \\
11 & Release audit, synchronized proof ledger and human-readable PDF &
High \\
\end{longtable}

Execute one prompt per milestone. Each produces a build log, declaration
inventory, assumption/axiom audit, updated proof ledger, and a report
against this architecture. Stop at the milestone boundary and submit the
result for review. Helper lemmas and API choices belong to Codex;
changes to the mathematical assumptions, scope, or conclusion require an
explicit design review. A blocker report gives the precise failed lemma
and attempted implementation route, not a request to let Codex weaken
the economics.

The supplied manifest contains 52 core contracts, two source-diagnostic
contracts, and three exact-extension contracts. A core release also
requires the diagnostics before it labels the source correction
certified.

Statuses are \texttt{UNFORMALIZED}, \texttt{IN\_PROGRESS},
\texttt{KERNEL\_CHECKED}, \texttt{REVIEW\_READY}, \texttt{GREEN}, and
\texttt{BLOCKED}. \texttt{GREEN} requires both kernel checking and an
accepted adequacy review. The final core release has no amber or
closure-assumed theorem. An excluded or deferred source claim is visible
in the coverage manifest; it is never reclassified as completed just
because the core builds.

\section{15. References and API
locators}\label{references-and-api-locators}

Aiyagari, S. Rao (1993). \emph{Uninsured Idiosyncratic Risk and
Aggregate Saving}. Federal Reserve Bank of Minneapolis Working Paper
502, revised December 1993.

Aiyagari, S. Rao (1994). ``Uninsured Idiosyncratic Risk and Aggregate
Saving.'' \emph{Quarterly Journal of Economics} 109(3), 659--684.

Benveniste, Lawrence M., and José A. Scheinkman (1979). ``On the
Differentiability of the Value Function in Dynamic Models of
Economics.'' \emph{Econometrica} 47(3), 727--732.

Chamberlain, Gary, and Charles A. Wilson (2000). ``Optimal Intertemporal
Consumption under Uncertainty.'' \emph{Review of Economic Dynamics}
3(3), 365--395. Later published source, not silently equated with the
1984 working paper cited by Aiyagari.

Clarida, Richard H. (1987). ``Consumption, Liquidity Constraints and
Asset Accumulation in the Presence of Random Income Fluctuations.''
\emph{International Economic Review} 28(2), 339--351.

Clarida, Richard H. (1990). ``International Lending and Borrowing in a
Stochastic, Stationary Equilibrium.'' \emph{International Economic
Review} 31(3), 543--558.

Miller, Bruce L. (1976). ``The Effect on Optimal Consumption of
Increased Uncertainty in Labor Income in the Multiperiod Case.''
\emph{Journal of Economic Theory} 13, 154--167.

Schechtman, Jack, and Vera L. S. Escudero (1977). ``Some Results on `An
Income Fluctuation Problem'.'' \emph{Journal of Economic Theory} 16(2),
151--166.

Sibley, David S. (1975). ``Permanent and Transitory Income Effects in a
Model of Optimal Consumption with Wage Income Uncertainty.''
\emph{Journal of Economic Theory} 11, 68--82.

Stokey, Nancy L., and Robert E. Lucas, Jr., with Edward C. Prescott
(1989). \emph{Recursive Methods in Economic Dynamics}. Harvard
University Press.

Official Mathlib and Codex API locators, inspected on September 11,
2026, are recorded in \texttt{contracts/api\_sources.json}. Prompt 00
checks their exact signatures against the pinned checkout; no installed
API or Lean build is presumed by this document.

\end{document}
